\section{Introduction}

We aim to understand the stochastics of those infections which exploit a lytic
mechanism in their initial pathogenesis: infectious agents that use intracellular
environments to their proliferative advantage. The importance of stochasticity
here can hardly be overstated. During the first handful of replicative cycles
there is an exaggerated chance that an infection dies out before it gets started,
and a deterministic treatment risks glossing over this perilous beginning, its
variables venturing smoothly between one bacterium and none. To grasp
\textit{why} \textit{Y.~pestis}, and other pathogens like it, do what they do,
we must account for the interplay of chance and fortune in the smallest
populations. It is for this reason that we have selected the
birth--death--catastrophe process as our tool of inquiry.

Study of the birth--death--catastrophe (BDC) process and its dynamics forms a
central focus of this thesis. Here we present a series of results and analyses
associated with its behaviour. As a continuous-time, discrete-state Markov
chain, the BDC process found its origin in a paper by Brockwell, Gani and
Resnick published in 1982 \cite{brockwell1982birth}. During the late 1970s and
early 1980s, attempts were being made to describe unpredictable population
collapses using stochastic process theory:
\begin{quotation}
``There has recently been a great deal of interest in the development and
analysis of stochastic models for the growth of populations which are subject to
catastrophes due either to death or large-scale emigrations''
\cite{brockwell1982birth}.
\end{quotation}
Prior to that work, models had largely incorporated deterministic growth
interspersed by stochastic reductions --- sudden drops, the sizes and times of
which were defined in distribution (see, for example, the models of Pakes and
collaborators from the late 1970s). These suffered from the peculiarities of
modelling discrete populations with continuous random variables, and so attempts
were made to employ a step process instead: a continuous-time Markov chain in a
discrete state space.

In fact, the model presented in that 1982 paper is a generalisation of what we
refer to as the BDC process. Not only did it also include immigration events ---
single steps up occurring as a Poisson process with constant rate --- but the
catastrophic steps down in population count were sized as random i.i.d.\
variables, with a few basic distributions investigated. In our case, the
catastrophe event reduces the population directly to $0$, or to a special
post-catastrophe state $H$; in either case causing fixation.

Our interest in the birth--death--catastrophe process stems from within-host
infection modelling. Like others before us, we are interested
in the pathogenic behaviour of intracellular replication and bursting. Our agent
of focus, \textit{Y.~pestis}, finds its way into host macrophages and replicates
there, before ultimately lysing the cell so that the population of infectious
bacteria can traverse the extracellular matrix and continue its exploitative
fight for survival. ``Catastrophe'' here names the cell-bursting event, which
becomes increasingly likely as the intracellular population expands. Carruthers
and collaborators worked on the same model with focus on the agent
\textit{Francisella tularensis} \cite{carruthers2020stochastic}, and similar
work has been directed towards the pathogenesis of salmonella. The mathematical
backbone of these models rests on the classical theory of birth-and-death
processes and their killed variants
\cite{karlin1957classification,karlin1982linear,di2008note}.

Our particular interest is the potential importance of this replication and
bursting strategy, along with other behaviours of \textit{Y.~pestis}, for an
advantage in early pathogenesis --- getting a foothold when numbers are few,
before a serious immune response arrives. Those questions return in
Chapters~4a and~4b, where the single-cell theory developed here is completed and
then embedded in a population model; for the purposes of the present chapter, we
take the birth--death--catastrophe process as an abstract model, unbound by any
specific application. We will present a number of mathematical results which
quantify the stochastics of its growth and fixation over time.

\subsection*{Further interest}

Interest in the dynamics of catastrophic population decline is nothing new.
Towards the end of the 18th century, Thomas Malthus published his fears of
impending societal collapse. He worried that a geometrically growing census
would be less and less supported by only arithmetic development in food
production, ultimately resulting in mass famine. In the relevant literature,
such population reductions by external factors are termed ``positive checks'',
as opposed to the ``preventative checks'' by which individuals have fewer
children to account for resource scarcity.

In 1755, just 11 years before Thomas Malthus was born, Lisbon experienced a
substantial earthquake that devastated the city. Voltaire includes a graphic
scene of this in his allegorical novel \textit{Candide}. Perhaps catastrophe was
on the mind of the age, and maybe the anxiety was justified. France, after all,
was on the brink of revolution:

\begin{quotation}
``War had created such social systems of unprecedented complexity that they
exceeded the capacities of their internal regulatory mechanisms to orchestrate
the behaviours of their parts. When, as in France, administrative chaos became
too great and the government collapsed --- an event easily measured by the
financial crisis of 1788 --- society exploded.'' \cite{artigiani1987revolution}
\end{quotation}

This quote is from a paper which attempts to apply Prigogine's theory of
dissipative structures to the evolution of culture and societal norms through
revolutions. It is perhaps ambitious, and lacking in empirical merit, but it has
a certain appeal. Similarly, we are curious about the prospect that the BDC
process, or a variant of it, can be taken in allegory to illustrate a general
behaviour observed in dissipative structures present across multiple dynamic
phenomena.

We have also been interested in the role of catastrophic events in triggering
evolutionary transitions. A notable example is the KT boundary, where
65 million years ago much of life on earth was wiped out by an asteroid,
setting the stage for a mass proliferation and diversification of mammals.

It has also been postulated that the catastrophes of plague in Europe and their
aftermath contributed to advancements in technology and culture --- most recently
by James Belich, who argues that previous estimates for the scale of population
decline sit below the real figure, claiming the initial wave saw a death of at
least 50\% of the European population in the decade following 1345. With such a
sudden elimination of so many people, labour was in reduced supply. It has been
hypothesised that a subsequent necessity to do more work with fewer people
encouraged a number of mechanical inventions; these, along with higher wages,
might have contributed to the improved prosperity and the cultural advancements
seen in the renaissance and the enlightenment periods. It is even said that the
Lisbon earthquake and its cleanup spurred advancements in geographical
techniques: damage to buildings was diligently recorded by locality, providing
new light on the concept of an epicentre.

These evolutionary echoes of catastrophe will be touched on again when, in
Chapter~4b, we compare the bursting strategy against its alternatives. For now,
let us take a step back, and then forward, into the mathematics of the
birth--death--catastrophe process.
